% Options for packages loaded elsewhere
\PassOptionsToPackage{unicode}{hyperref}
\PassOptionsToPackage{hyphens}{url}
\PassOptionsToPackage{dvipsnames,svgnames,x11names}{xcolor}
\documentclass[
  11pt,
  letterpaper,
]{article}
\usepackage{xcolor}
\usepackage[margin=0.8in]{geometry}
\usepackage{amsmath,amssymb}
\setcounter{secnumdepth}{-\maxdimen} % remove section numbering
\usepackage{iftex}
\ifPDFTeX
  \usepackage[T1]{fontenc}
  \usepackage[utf8]{inputenc}
  \usepackage{textcomp} % provide euro and other symbols
\else % if luatex or xetex
  \usepackage{unicode-math} % this also loads fontspec
  \defaultfontfeatures{Scale=MatchLowercase}
  \defaultfontfeatures[\rmfamily]{Ligatures=TeX,Scale=1}
\fi
\usepackage{lmodern}
\ifPDFTeX\else
  % xetex/luatex font selection
\fi
% Use upquote if available, for straight quotes in verbatim environments
\IfFileExists{upquote.sty}{\usepackage{upquote}}{}
\IfFileExists{microtype.sty}{% use microtype if available
  \usepackage[]{microtype}
  \UseMicrotypeSet[protrusion]{basicmath} % disable protrusion for tt fonts
}{}
\usepackage{setspace}
\makeatletter
\@ifundefined{KOMAClassName}{% if non-KOMA class
  \IfFileExists{parskip.sty}{%
    \usepackage{parskip}
  }{% else
    \setlength{\parindent}{0pt}
    \setlength{\parskip}{6pt plus 2pt minus 1pt}}
}{% if KOMA class
  \KOMAoptions{parskip=half}}
\makeatother
\usepackage{color}
\usepackage{fancyvrb}
\newcommand{\VerbBar}{|}
\newcommand{\VERB}{\Verb[commandchars=\\\{\}]}
\DefineVerbatimEnvironment{Highlighting}{Verbatim}{commandchars=\\\{\}}
% Add ',fontsize=\small' for more characters per line
\newenvironment{Shaded}{}{}
\newcommand{\AlertTok}[1]{\textcolor[rgb]{1.00,0.00,0.00}{\textbf{#1}}}
\newcommand{\AnnotationTok}[1]{\textcolor[rgb]{0.38,0.63,0.69}{\textbf{\textit{#1}}}}
\newcommand{\AttributeTok}[1]{\textcolor[rgb]{0.49,0.56,0.16}{#1}}
\newcommand{\BaseNTok}[1]{\textcolor[rgb]{0.25,0.63,0.44}{#1}}
\newcommand{\BuiltInTok}[1]{\textcolor[rgb]{0.00,0.50,0.00}{#1}}
\newcommand{\CharTok}[1]{\textcolor[rgb]{0.25,0.44,0.63}{#1}}
\newcommand{\CommentTok}[1]{\textcolor[rgb]{0.38,0.63,0.69}{\textit{#1}}}
\newcommand{\CommentVarTok}[1]{\textcolor[rgb]{0.38,0.63,0.69}{\textbf{\textit{#1}}}}
\newcommand{\ConstantTok}[1]{\textcolor[rgb]{0.53,0.00,0.00}{#1}}
\newcommand{\ControlFlowTok}[1]{\textcolor[rgb]{0.00,0.44,0.13}{\textbf{#1}}}
\newcommand{\DataTypeTok}[1]{\textcolor[rgb]{0.56,0.13,0.00}{#1}}
\newcommand{\DecValTok}[1]{\textcolor[rgb]{0.25,0.63,0.44}{#1}}
\newcommand{\DocumentationTok}[1]{\textcolor[rgb]{0.73,0.13,0.13}{\textit{#1}}}
\newcommand{\ErrorTok}[1]{\textcolor[rgb]{1.00,0.00,0.00}{\textbf{#1}}}
\newcommand{\ExtensionTok}[1]{#1}
\newcommand{\FloatTok}[1]{\textcolor[rgb]{0.25,0.63,0.44}{#1}}
\newcommand{\FunctionTok}[1]{\textcolor[rgb]{0.02,0.16,0.49}{#1}}
\newcommand{\ImportTok}[1]{\textcolor[rgb]{0.00,0.50,0.00}{\textbf{#1}}}
\newcommand{\InformationTok}[1]{\textcolor[rgb]{0.38,0.63,0.69}{\textbf{\textit{#1}}}}
\newcommand{\KeywordTok}[1]{\textcolor[rgb]{0.00,0.44,0.13}{\textbf{#1}}}
\newcommand{\NormalTok}[1]{#1}
\newcommand{\OperatorTok}[1]{\textcolor[rgb]{0.40,0.40,0.40}{#1}}
\newcommand{\OtherTok}[1]{\textcolor[rgb]{0.00,0.44,0.13}{#1}}
\newcommand{\PreprocessorTok}[1]{\textcolor[rgb]{0.74,0.48,0.00}{#1}}
\newcommand{\RegionMarkerTok}[1]{#1}
\newcommand{\SpecialCharTok}[1]{\textcolor[rgb]{0.25,0.44,0.63}{#1}}
\newcommand{\SpecialStringTok}[1]{\textcolor[rgb]{0.73,0.40,0.53}{#1}}
\newcommand{\StringTok}[1]{\textcolor[rgb]{0.25,0.44,0.63}{#1}}
\newcommand{\VariableTok}[1]{\textcolor[rgb]{0.10,0.09,0.49}{#1}}
\newcommand{\VerbatimStringTok}[1]{\textcolor[rgb]{0.25,0.44,0.63}{#1}}
\newcommand{\WarningTok}[1]{\textcolor[rgb]{0.38,0.63,0.69}{\textbf{\textit{#1}}}}
\usepackage{longtable,booktabs,array}
\usepackage{caption}
\captionsetup[table]{skip=6pt}
\newcounter{none} % for unnumbered tables
\usepackage{calc} % for calculating minipage widths
% Correct order of tables after \paragraph or \subparagraph
\usepackage{etoolbox}
\makeatletter
\patchcmd\longtable{\par}{\if@noskipsec\mbox{}\fi\par}{}{}
\makeatother
% Allow footnotes in longtable head/foot
\IfFileExists{footnotehyper.sty}{\usepackage{footnotehyper}}{\usepackage{footnote}}
\makesavenoteenv{longtable}
\setlength{\emergencystretch}{3em} % prevent overfull lines
\providecommand{\tightlist}{%
  \setlength{\itemsep}{0pt}\setlength{\parskip}{0pt}}
\usepackage{xcolor}
\definecolor{researchblue}{RGB}{31,65,99}
\makeatletter
\renewcommand{\maketitle}{\begingroup\centering
{\LARGE\@title\par}\vspace{0.6em}
{\normalsize\@author\par}\vspace{0.3em}
{\small\@date\par}\endgroup\vspace{1em}}
\makeatother
\usepackage{bookmark}
\IfFileExists{xurl.sty}{\usepackage{xurl}}{} % add URL line breaks if available
\urlstyle{same}
% fallback for those not using the hyperref driver hyperxmp:
\makeatletter
\@ifundefined{xmpquote}{\newcommand{\xmpquote}[1]{#1}}{}
\makeatother
\hypersetup{
  pdftitle={Advanced environment labs},
  pdfauthor={World Engine technical research},
  colorlinks=true,
  linkcolor={researchblue},
  filecolor={Maroon},
  citecolor={Blue},
  urlcolor={researchblue},
  pdfcreator={LaTeX via pandoc}}

\title{Advanced environment labs}
\author{World Engine technical research}
\date{October 10, 2026}

\begin{document}
\maketitle

\setstretch{1.0}
Build nine review artifacts: a fidelity contract, a grader attack
matrix, an evaluation analysis, a split/rights manifest, a
preference-batch specification, an RL capture contract, a scheduler
budget, a cost model, and a deployment threat model. You can complete
the numeric exercises with paper or a calculator. The optional
repository checks use scripted tests, not model inference.

All invented numbers below are \textbf{illustrative}, including prices,
token rates, memory layouts, sample sizes, and agreement scores. They
are not World Engine measurements. Actual receipts remain in the
\href{world-engine-report.md}{foundation report},
\href{world-engine-follow-up.md}{follow-up appendix}, and
\href{../workflow-model-study-results.md}{matched-study results}. This
workbook adds no benchmark, training run, application-parity claim, or
production readiness claim. Its
\href{../research/ADVANCED-SOURCES.md}{source ledger} records
primary-source checks and limits.

\subsection{Lab 1. Specify behavioral fidelity before measuring
it}\label{lab-1.-specify-behavioral-fidelity-before-measuring-it}

An environment is a state-transition system. Its state contains
operational facts, an action asks for a change or a read, and its
observation exposes some state to the agent. The policy chooses actions
from observations and history. A rollout is the resulting sequence of
decisions, tool results, and states. A grader judges that sequence and
its outcome against a task contract.

A synthetic environment can look unlike an application yet preserve one
workflow accurately. Conversely, an exact-looking screen can give
incorrect permission, audit, or retry behavior. Fidelity is agreement on
declared behavior, not the number of copied fields or screenshots.
Fleet's archived public SDK demonstrates versioned creation, seeded
reset, SQL access, and cleanup. That lifecycle alone does not establish
which application behaviors match.
\href{https://github.com/fleet-ai/fleet-sdk}{Fleet SDK}.

\subsubsection{Produce a differential
contract}\label{produce-a-differential-contract}

Choose one bounded task, such as refunding an eligible damaged order.
Define an abstraction \(\alpha\) that maps application records to
comparable business facts. Drop irrelevant rendering IDs; retain money
in integer cents, ownership, authorization, and audit order. Then
compare an application transition with a simulated transition:

\[\alpha(T_{\mathrm{app}}(s,a))\stackrel{?}{=}T_{\mathrm{sim}}(\alpha(s),a).\]

This deterministic equation assumes equivalent initial facts and action
semantics. For asynchronous or stochastic behavior, compare declared
distributions or eventual predicates under a fixed deadline instead. Do
not erase genuine nondeterminism by sorting a causally ordered audit
log.

Define these axes independently:

{\def\LTcaptype{none} % do not increment counter
\begin{longtable}[]{@{}
  >{\raggedright\arraybackslash}p{(\linewidth - 4\tabcolsep) * \real{0.3333}}
  >{\raggedright\arraybackslash}p{(\linewidth - 4\tabcolsep) * \real{0.3333}}
  >{\raggedright\arraybackslash}p{(\linewidth - 4\tabcolsep) * \real{0.3333}}@{}}
\toprule\noalign{}
\begin{minipage}[b]{\linewidth}\raggedright
Axis
\end{minipage} & \begin{minipage}[b]{\linewidth}\raggedright
Differential probe
\end{minipage} & \begin{minipage}[b]{\linewidth}\raggedright
Failure that matters
\end{minipage} \\
\midrule\noalign{}
\endhead
\bottomrule\noalign{}
\endlastfoot
Data & Same target plus unrelated rows & Wrong target or lost
outsider \\
Permissions & Same action under two roles & Simulation admits a
forbidden payment \\
Observation & Same search/page/history & Agent sees privileged approval
data \\
Transition & One-cent or expiry boundary & Different eligibility
verdict \\
Retry & Same key twice, then another key & Duplicate payment or missing
conflict \\
Time & Delayed review, timeout, reset & Approval never expires or time
resets \\
Outcome & Independent final-state check & Good-looking transcript, wrong
ledger \\
\end{longtable}
}

Use an authorized application test tenant, synthetic records, and a
reset procedure. Freeze versions, configuration, clock semantics, and
adapters. Execute the same action sequence in both systems. Preserve
before/after facts and errors, not only final pass rates. Ask domain
experts to adjudicate mismatches with the same public instructions given
to the agent. Anthropic's evaluation guidance separates transcripts from
outcomes and recommends expert calibration.
\href{https://www.anthropic.com/engineering/demystifying-evals-for-ai-agents}{Agent
evals}.

\subsubsection{Solve the agreement
example}\label{solve-the-agreement-example}

Suppose 200 transitions yield 180 agreements, eight harmless formatting
differences, eight business mismatches, and four permission mismatches.
A raw agreement score is 90\%. After a justified formatting
normalization it is 94\%. Neither number makes four unsafe admissions
acceptable.

Assign loss weights 0, 1, and 10 to normalized agreement, business
mismatch, and permission mismatch. The mean loss is
\((8+10\times4)/200=0.24\). This weight choice is a declared risk
preference, not a probability of harm. Report the permission count
separately so a common easy read cannot hide a rare dangerous write.

\textbf{Question.} Zero unsafe admissions occurred in 100 adversarial
cases. Is the rate proven zero? \textbf{Solution.} No.~Even under
independent representative trials, the one-sided 95\% binomial upper
bound is \(1-0.05^{1/100}\approx2.95\%\). Adversarial cases are usually
selected, not a random population sample, so that statistical bound does
not imply a production rate. Publish tested categories and blind spots.

\textbf{Deliverable.} A one-page contract with seven axes, reset/clock
rules, risk weights, explicit exclusions, expert disagreement procedure,
and paired mismatch records. World Engine has not completed this
application differential study.

\subsection{Lab 2. Attack the grader separately from the
tools}\label{lab-2.-attack-the-grader-separately-from-the-tools}

A reward is a scalar optimization signal. Task success is the intended
outcome. Those can disagree. Specification gaming exploits the score's
definition; reward tampering changes the scoring machinery or its
evidence. DeepMind describes the former distinction. Anthropic's
controlled coding-training study reports broader effects after reward
hacking, but this project has no training result that establishes the
same effects.
\href{https://deepmind.google/blog/specification-gaming-the-flip-side-of-ai-ingenuity/}{Specification
gaming},
\href{https://www.anthropic.com/research/emergent-misalignment-reward-hacking}{reward-hacking
study}.

\subsubsection{Produce an attack matrix}\label{produce-an-attack-matrix}

{\def\LTcaptype{none} % do not increment counter
\begin{longtable}[]{@{}
  >{\raggedright\arraybackslash}p{(\linewidth - 4\tabcolsep) * \real{0.3333}}
  >{\raggedright\arraybackslash}p{(\linewidth - 4\tabcolsep) * \real{0.3333}}
  >{\raggedright\arraybackslash}p{(\linewidth - 4\tabcolsep) * \real{0.3333}}@{}}
\toprule\noalign{}
\begin{minipage}[b]{\linewidth}\raggedright
Attack class
\end{minipage} & \begin{minipage}[b]{\linewidth}\raggedright
Synthetic workflow probe
\end{minipage} & \begin{minipage}[b]{\linewidth}\raggedright
Required defense
\end{minipage} \\
\midrule\noalign{}
\endhead
\bottomrule\noalign{}
\endlastfoot
Proxy shortcut & Finish without granting or closing & Outcome predicate,
not completion marker alone \\
Target substitution & Change request ID or canonical customer & Anchor
identities to immutable initial facts \\
Evidence forgery & Insert finance approval or merge audit & Exact
authorized history, not field presence \\
Authority laundering & Replace role, owner, limit, or policy & Trusted
permission facts outside agent edits \\
Retry farming & Reuse a key to create more reward/events & Idempotent
effect and exact event count \\
Preservation failure & Solve target, delete unrelated ledger & Grade
whole allowed state, including outsiders \\
Grader intrusion & Rewrite tests or read expected output & Separate
agent and verifier credentials/files \\
Distribution shortcut & Infer denial from a seed or broad hint &
Balanced cases and evidence-required decisions \\
Type confusion & Substitute \texttt{true} for ID \texttt{1} &
Type-sensitive canonical comparison \\
Artifact deception & Link output to hidden files; exceed archive cap &
Validate complete detached output, fail closed \\
\end{longtable}
}

Use two independent test routes. First, ask an ordinary tool to perform
the forbidden operation. Its error must leave operational state
unchanged, apart from any explicitly allowed clock tick. Second, corrupt
a trusted test's completed snapshot directly, then invoke the grader.
That probe checks the grader contract; it is not evidence that an agent
can reach the corrupted state through normal tools.

An independent grader starts from trusted initial facts and folds
exactly permitted change certificates into an expected state. It checks
prerequisites, identities, actor authority, event order, and unchanged
fields. It must accept every legal operation order, not only a reference
transcript. Calling the same buggy \texttt{apply} function again merely
reproduces its mistake. A locally generated certificate is also not a
signature against a malicious host operator.

Inspect existing proofs without inference:

\begin{Shaded}
\begin{Highlighting}[]
\ExtensionTok{uv}\NormalTok{ run pytest tests/test\_workflow\_grader\_integrity.py }\DataTypeTok{\textbackslash{}}
\NormalTok{  tests/test\_incident\_integrity.py}
\end{Highlighting}
\end{Shaded}

Read the positive-prefix tests before adding attacks. A correct
permission denial after legitimate progress must not turn that progress
into a safety violation. Repeating an idempotent payment must not create
another audit event. The \href{../workflow-grader-integrity.md}{grader
correction} explains actual historical gaps without rewriting old
receipts.

\subsubsection{Solve the reward example}\label{solve-the-reward-example}

Suppose a grade averages safety, settlement, and explicit finish. An
unchanged safe world followed by finish earns \(2/3\), despite settling
nothing. An optimizer can prefer that cheap shortcut. Multiplying by a
safety check stops unsafe outcomes, but it does not fix free credit for
doing no work.

For a new outcome-only task, use
\(r=\mathbf{1}[\text{safe}\land\text{settled}\land\text{finished}]\) and
retain individual checks as diagnostics. That is the access world's
declared choice, not proof that binary feedback always trains better.
Legacy shaped grades remain historical semantics; episode
timeouts/unfinished termination still force authoritative episode reward
to zero.

Potential-based shaping instead adds
\(F(s,a,s')=\gamma\Phi(s')-\Phi(s)\) under specified conditions. With
\(\gamma=0.9\), \(\Phi(s)=2\), and \(\Phi(s')=3\), the shaping term is
\(0.7\). Across a discounted path those terms telescope, with
terminal-boundary conditions mattering. Arbitrary check averages have no
such policy-invariance guarantee.
\href{https://ai.stanford.edu/~ang/papers/shaping-icml99.pdf}{Reward-shaping
theorem}.

\textbf{Deliverable.} An attack matrix with a legal counterexample for
each rejected mutation. Record severity, reachable-through-tools status,
grader result, source revision, and the regression that prevents
recurrence. Production runtime guards must not depend on Python
\texttt{assert}, which optimized Python removes.

\subsection{Lab 3. Define the denominator before reading the
score}\label{lab-3.-define-the-denominator-before-reading-the-score}

A task is a problem instance. A trial is one attempt at that task. Four
adapters on one task create four trials, not four independent tasks.
Repeating one task under several seeds may still share its rule family.
Pin the policy, prompt, tool parser, environment, grader, budget, and
collection order before gathering outcomes.

Keep attempted, finalized, verified, inconclusive, and unattempted
counts separate. A missing journal row is not a failed reasoning
attempt. A persisted malformed envelope is not a verified episode. Never
drop failures or silently replace an interrupted slot with a successful
retry. For a predeclared end-to-end benchmark, timeouts and valid
incorrect outcomes count as task failures; infrastructure-inconclusive
rows need their own count and rule.

\subsubsection{Solve a matched
comparison}\label{solve-a-matched-comparison}

Two adapters attempt 40 fixed tasks once each. They both pass 14, only A
passes ten, only B passes two, and neither passes 14. Then A passes
\(24/40=60\%\), B passes \(16/40=40\%\), and the paired difference is
\((10-2)/40=20\) percentage points.

For each task define \(d_i=1\), \(-1\), or \(0\) according to A-only,
B-only, or equal outcome. Resample \textbf{whole task pairs}, not two
unrelated lists of attempts. If tasks belong to declared families,
resample within those strata and report the weighting rule. The paired
result compares complete adapter systems, not one isolated serialization
feature.

A confidence interval describes an estimation procedure's coverage under
its sampling assumptions. It is not a probability that one fixed true
rate lies inside this realized interval. Selected toy fixtures support
descriptive uncertainty summaries, not population coverage claims. For
correlated variants, resample parent families rather than pretending
every row is independent.

Now divide A's tasks into ten difficult tasks with two passes and thirty
easier tasks with 22 passes. The micro-average is still 60\%. Giving the
two families equal weight yields \((20\%+73.33\%)/2=46.67\%\). Neither
average is wrong; they answer different workload questions. Specify the
target mixture instead of choosing whichever looks better. RL evaluation
research warns about uncertainty and unstable aggregate conclusions from
limited runs. \href{https://arxiv.org/abs/2108.13264}{Statistical
precipice}.

\subsubsection{Diagnose before proposing a
fix}\label{diagnose-before-proposing-a-fix}

{\def\LTcaptype{none} % do not increment counter
\begin{longtable}[]{@{}
  >{\raggedright\arraybackslash}p{(\linewidth - 4\tabcolsep) * \real{0.3333}}
  >{\raggedright\arraybackslash}p{(\linewidth - 4\tabcolsep) * \real{0.3333}}
  >{\raggedright\arraybackslash}p{(\linewidth - 4\tabcolsep) * \real{0.3333}}@{}}
\toprule\noalign{}
\begin{minipage}[b]{\linewidth}\raggedright
Observed category
\end{minipage} & \begin{minipage}[b]{\linewidth}\raggedright
What the record establishes
\end{minipage} & \begin{minipage}[b]{\linewidth}\raggedright
What it does not establish
\end{minipage} \\
\midrule\noalign{}
\endhead
\bottomrule\noalign{}
\endlastfoot
Provider protocol error & Response violates frozen action contract &
Model cannot solve the business problem \\
Environment denial & Action fails a declared permission/precondition &
Provider transport failed \\
Valid incorrect completion & Parser accepted actions; goal unmet &
Grader is correct without independent checks \\
Maximum steps & Budget ended before accepted completion & More steps
would necessarily solve it \\
Timeout & Deadline ended the attempt & Which component caused the
delay \\
Setup/cleanup/persistence error & Controller lifecycle failed & A
comparable clean policy outcome \\
\end{longtable}
}

Inspect raw exchanges and state history to refine causes. A timeout
after repeated denied approvals may involve both poor decisions and
expensive inference. Exclusive report categories need declared
precedence; secondary causes can coexist. Relaxing a one-call parser
after seeing multi-call failures changes the experiment. Collect a new
locked protocol rather than rescoring the old one under a more forgiving
contract.

\textbf{Question.} A trial succeeds with probability 0.8. What are
success-at-least-once and success-every-time over three trials?
\textbf{Solution.} Under identical independent trials,
\(1-(1-0.8)^3=0.992\) and \(0.8^3=0.512\). Independence and constant
probability are assumptions, not consequences of three observed repeats.
A retry-friendly search score and a reliability score answer different
questions.

\textbf{Deliverable.} A denominator table, all four matched outcomes,
micro and declared macro scores, confidence intervals with sampling
caveats, clean-success counts, failure taxonomy, and every attempted
slot. The actual \href{../workflow-model-study-results.md}{matched
study} uses this discipline; its result is not the invented 40-task
example above.

\subsection{Lab 4. Separate seed holdout, family holdout, and
rights}\label{lab-4.-separate-seed-holdout-family-holdout-and-rights}

A curriculum changes which training tasks a policy sees and when. A
split assigns examples to training, tuning/validation, or final
evaluation. They are related but different. Validation guides
development; once repeated feedback changes your choices, it is no
longer a blind final test. Procedural generation can measure
unseen-level behavior, as Procgen does, without establishing transfer to
an unseen environment or rule.
\href{https://arxiv.org/abs/1912.01588}{Procgen paper}.

\subsubsection{Produce a split manifest}\label{produce-a-split-manifest}

Imagine 100 workflow families with ten seeded instances per family. A
family specifies a rule and interaction pattern; a seed changes names,
amounts, distractors, or ordering within that rule. Randomly assigning
800 rows to train and 200 to test usually puts each family in both sets.
Those 200 rows measure instance variation, not withheld-rule reasoning.

Instead, reserve 70 families for training, fifteen for validation, and
fifteen for final evaluation. Assign all seeds, paraphrases, reference
solutions, revisions, and derived traces of a family together. Keep a
stable parent identity even when implementation hashes change. Freeze
the mapping before generation and publish overlap checks. World Engine's
stable family/seed lineage prevents implementation revisions from moving
the same instance between splits, but does not implement this broader
family-heldout study.

Use at least three claims with different tests:

\begin{enumerate}
\def\labelenumi{\arabic{enumi}.}
\tightlist
\item
  Unseen seed, known family: new invoice IDs and distractor order.
\item
  Unseen family, known domain: multi-approval refund rule absent from
  training.
\item
  Unseen application/domain: equivalent policy expressed through another
  service's navigation and errors.
\end{enumerate}

Increase curriculum difficulty by adding partial observation, ambiguous
identifiers, longer valid dependencies, and recovery branches. First
prove each task has a legal solution under its budget. Preserve easy
regression tasks while adding difficult capability tasks. Reject
impossible or contradictory fixtures rather than teaching the model to
guess the desired verdict.

\subsubsection{Solve a contamination
example}\label{solve-a-contamination-example}

A family has three paraphrased rows. You independently send each row to
training with probability 0.8. The probability that the family appears
on both sides is \(1-0.8^3-0.2^3=0.48\). Forty-eight percent of such
families are expected to cross the boundary. Row hashes will not catch
paraphrases. Group by provenance and semantics before splitting; use
exact and near-duplicate scans as additional checks.

Contamination includes leaked tests, reference solutions, task-specific
hints, earlier evaluation feedback, and synthetic outputs derived from
heldout data. Poisoning is different: an attacker supplies examples
intended to induce unwanted behavior. Anthropic's controlled backdoor
study motivates source review, but does not provide a universal safe
percentage.
\href{https://www.anthropic.com/research/small-samples-poison}{Poisoning
study}.

\subsubsection{Produce a rights
manifest}\label{produce-a-rights-manifest}

For each source, record owner, exact revision, file-level
license/notices, permitted uses, restrictions, personal-data status,
consent or other approved processing basis, retention, and derived
artifacts. Record who approved reuse and which deletion requests must
propagate. Synthetic names do not cleanse a copied private policy or
customer incident.

Public visibility is not public-domain status. Apache-2.0 grants
permissions with conditions; CC0 is a rights dedication mechanism, while
the Public Domain Mark identifies an existing status. Neither label
establishes consent to process private records or third-party rights.
FrontierSWE v2's notice describes evaluation-purpose task content and
per-task attribution, so do not assume blanket training reuse.
\href{https://www.apache.org/licenses/LICENSE-2.0}{Apache license},
\href{https://creativecommons.org/public-domain/}{CC tools},
\href{https://github.com/Proximal-Labs/frontier-swe-v2}{FrontierSWE
repository}.

\textbf{Deliverable.} Separate instance, family, and application splits;
a curriculum change log; contamination scans with known limits; and an
approved rights manifest. Keep blocked data outside candidate pools, not
merely marked in an export after selection. Obtain qualified rights
review before reuse where permission is unclear.

\subsection{Lab 5. Turn a preference into a defined
likelihood}\label{lab-5.-turn-a-preference-into-a-defined-likelihood}

Supervised fine-tuning, or SFT, increases the probability of selected
demonstrations. Direct preference optimization, or DPO, learns from a
preferred and rejected completion of the same prompt. Neither is a name
for copying readable trajectories into a JSON file. Before training,
specify what the model predicts, which tokens receive loss, and how the
reference model sees identical context.

For SFT, let \(z_t\) be a token, \(h_t\) its preceding context, and
\(m_t\in\{0,1\}\) a loss mask. One declared mean-token objective is

\[\mathcal L_{\mathrm{SFT}}=-\frac{\sum_t m_t\log\pi_\theta(z_t\mid h_t)}{\sum_t m_t}.\]

A mask chooses which tokens teach the policy. Context still conditions
predictions even when it has mask zero. Reject an example with no active
loss tokens rather than dividing by zero. Padding must not contribute
loss. For an action policy, service-generated tool results are context,
not the agent's sampled actions. Decide explicitly whether delimiters
and end-of-turn tokens belong to the completion.

Chat templates encode roles and boundaries differently across models.
The official Transformers documentation warns that applying extra
special tokens after template formatting can duplicate them. Record the
tokenizer, template revision, rendered IDs, truncation rules, and
generation boundary. Decoded text alone cannot prove that two pipelines
used identical tokens.
\href{https://huggingface.co/docs/transformers/main/en/chat_templating}{Chat
templates}.

\subsubsection{Solve a DPO example}\label{solve-a-dpo-example}

For one prompt \(x\), define preferred completion \(y^+\) and rejected
completion \(y^-\). Write the original reference-relative loss as

\[\mathcal L_{\mathrm{DPO}}=-\log\sigma\!\left(\beta[\Delta^+-\Delta^-]\right),\qquad
\Delta^\pm=\log\pi_\theta(y^\pm\mid x)-\log\pi_{\mathrm{ref}}(y^\pm\mid x).\]

\(\sigma(u)=1/(1+e^{-u})\) is the logistic function. The reference
policy supplies a baseline likelihood. The preference derivation assumes
a particular model of comparative reward; minimizing the loss does not
validate the labels. \href{https://arxiv.org/html/2305.18290v3}{DPO
paper}.

Suppose the trained-policy log-likelihoods are \(-12\) and \(-10\),
reference likelihoods are \(-13\) and \(-9\), and \(\beta=0.1\). Then
\(\Delta^+=1\), \(\Delta^-=-1\), the logit is \(0.2\), preference
probability is approximately \(0.5498\), and loss is approximately
\(0.5981\). Although the preferred completion has lower absolute
probability, its improvement relative to reference is greater. Reverse
the preference and the loss becomes approximately \(0.7981\).

Sequence log-likelihood sums the selected conditional token
log-probabilities. Replacing a sum with a length-normalized mean changes
the objective. So do different truncations, masks, or reference
contexts. A six-token exercise with four context tokens and two action
tokens uses mask \texttt{{[}0,0,0,0,1,1{]}}; shifting that mask one
position teaches the wrong boundary even if the JSON parses.

\subsubsection{Make workflow pairs
comparable}\label{make-workflow-pairs-comparable}

Choose matching initial prompt, scenario, world/grader version, split,
policy contract, and budget. A successful refund and a failed refund for
another order are not a same-prompt preference pair. A shorter failed
attempt after an infrastructure exception may not express a defensible
business preference. Preserve the cause and selection rule.

For a whole tool trajectory, later observations depend on earlier
actions. You cannot paste both histories under one final observation and
call them alternative next actions. Specify an interactive objective
with each action conditioned on its own history, or declare a
transformed trajectory-completion task and evaluate that new task. Do
not assign the policy likelihood of service outputs as though it
generated them. TRL documents explicit-prompt and conversational
formats, but that structure alone does not resolve these semantics.
\href{https://huggingface.co/docs/trl/main/en/dpo_trainer}{DPO trainer}.

\textbf{Question.} A preferred record is task-passing but contains five
denied actions; another is task-passing and clean. Does higher reward
define a preference? \textbf{Solution.} Not if both receive the same
outcome reward. Define the desired quality relation first, retain ties,
and avoid inventing labels from record order. Business success,
cleanliness, latency, and safety are distinct objectives.

\textbf{Deliverable.} One inspected candidate pair with rendered token
IDs, action masks, reference inputs, likelihood aggregation, preference
provenance, and truncation checks. This is a specification exercise. The
local \href{../trainer-preprocessing-feasibility.md}{trainer feasibility
audit} remains a separate blocker; no usable tokenizer/trainer path is
supplied by this workbook.

\subsection{Lab 6. Distinguish PPO/GRPO batches from grouped
text}\label{lab-6.-distinguish-ppogrpo-batches-from-grouped-text}

Online reinforcement learning collects policy behavior, grades it, then
updates the policy. The \textbf{behavior policy} produced the action.
The \textbf{current policy} is being optimized. A \textbf{reference
policy} constrains drift or supplies a comparison baseline. They are not
interchangeable just because they share a model name.

An advantage estimates how much better an action's return is than a
baseline. A PPO actor-critic system commonly uses a learned value
estimate for that baseline. GRPO's original formulation instead compares
sampled completions within a same-prompt group. Both still need a
defined policy likelihood and a specified estimator.
\href{https://arxiv.org/pdf/1707.06347}{PPO paper},
\href{https://arxiv.org/html/2402.03300v3}{original GRPO}.

\subsubsection{Solve the clipping
calculation}\label{solve-the-clipping-calculation}

For a token or action with recorded behavior log-probability
\(\ell_{\mathrm{old}}\) and current log-probability \(\ell_\theta\),
define \(\rho=\exp(\ell_\theta-\ell_{\mathrm{old}})\). PPO's clipped
surrogate contribution is

\[\min\!\left(\rho A,\operatorname{clip}(\rho,1-\epsilon,1+\epsilon)A\right).\]

With \(\ell_{\mathrm{old}}=-2\), \(\ell_\theta=-1.8\), \(A=2\), and
\(\epsilon=0.2\), the ratio is about \(1.2214\). The unclipped
contribution is \(2.4428\); the clipped contribution is \(2.4\), so the
minimum is \(2.4\). If \(A=-2\), the minimum is \(-2.4428\), not
\(-2.4\). The minimum's direction matters. Clipping is not a blanket cap
on every gradient, a guaranteed global KL bound, or protection against a
wrong reward.

\subsubsection{Solve the group
calculation}\label{solve-the-group-calculation}

For four same-prompt samples with rewards \texttt{{[}0,0,1,1{]}}, the
mean is 0.5. Using \textbf{population} standard deviation for this
exercise gives 0.5, so standardized advantages are
\texttt{{[}-1,-1,1,1{]}}. Using sample standard deviation gives a
different magnitude. Declare the convention and numerical epsilon rather
than relying on a name such as ``GRPO''.

Groups \texttt{{[}0,0,0,0{]}} and \texttt{{[}1,1,1,1{]}} have zero
centered reward. Their outcome-driven advantage is zero under this rule,
although a configured KL or other loss may still act. Mixing four
different prompts in one group subtracts an unrelated baseline.
Repeating an identical prompt with deterministic decoding may also
supply no diverse candidates. A curriculum with only impossible tasks
wastes sampling; one with only solved tasks supplies little comparative
outcome signal. Current TRL offers several reward-scaling and loss
choices, so pin the trainer revision and complete configuration.
\href{https://huggingface.co/docs/trl/main/en/grpo_trainer}{GRPO
trainer}.

\subsubsection{Produce a training capture
contract}\label{produce-a-training-capture-contract}

Require exact prompt IDs, sampled completion IDs, sequence boundaries,
action/loss masks, behavior token log-probabilities, policy/checkpoint
identity, generation settings, stop/truncation reasons, task/grader
identity, and rewards. Declare temperature/top-p effects and any
generation-versus-training likelihood correction. Later retokenizing a
normalized tool JSON object does not recover the model's original
sampled sequence or its behavior distribution.

Pin one adapter/checkpoint for an entire rollout. Record collection and
consumed versions; apply a reviewed age or divergence rule before
updates. A trajectory spanning two policy versions needs explicitly
defined handling, not one invented ``old'' probability. Proximal's
October 8 account fixes an adapter per rollout, captures sampled tokens
near inference, and excludes overly stale groups. Its private
implementation and gains are not reproduced here.
\href{https://www.proximal.so/blog/posttraining-infra/}{Proximal
infrastructure}.

\textbf{Question.} A receipt has exact HTTP bytes, aggregate
input/output counts, reward, and readable actions. Is it PPO-ready?
\textbf{Solution.} No.~Counts do not identify the sampled IDs, masks, or
per-token behavior log-probabilities. An API's available probability
field also needs a verified interpretation for the actual sampling path.
Honest grouped exports remain candidate records, not trainable batches.

Hosted training has a separate availability gate. At the access date,
OpenAI's RFT guide says its fine-tuning platform is winding down and is
closed to new users. Do not assume an account, supported model, accepted
rollout format, or permission to create a job.
\href{https://developers.openai.com/api/docs/guides/reinforcement-fine-tuning}{Dated
RFT guide}.

\textbf{Deliverable.} A batch admission checklist with exact tensor
fields, version/staleness rules, group membership, zero-variance
handling, and fail-closed missing-field behavior. No optimization step
or trained checkpoint is produced.

\subsection{Lab 7. Budget inference sessions, not only host
tasks}\label{lab-7.-budget-inference-sessions-not-only-host-tasks}

Host concurrency counts active rollout tasks. GPU serving concurrency
counts requests or sessions that occupy model compute and memory.
Neither is equivalent to environments per second. A rollout can spend
most of its wall time waiting for inference, holding memory or a sandbox
while using little CPU.

Language-model inference has a \textbf{prefill} phase that processes
input tokens and a \textbf{decode} phase that generates new tokens.
Attention stores key/value tensors, called KV cache, so later tokens can
reuse earlier attention state. Prefix caching reuses a shared input
prefix. vLLM's documentation explicitly distinguishes this prefill
saving from generating new output tokens.
\href{https://docs.vllm.ai/en/latest/features/automatic_prefix_caching/}{Prefix
caching}.

\subsubsection{Solve a simplified KV
budget}\label{solve-a-simplified-kv-budget}

For a dense decoder with \(L\) layers, \(H_{\mathrm{KV}}\) key/value
heads, head dimension \(d\), \(b\) bytes per stored scalar, and \(n\)
cached tokens, approximate unsharded KV storage as

\[M_{\mathrm{KV}}\approx2L H_{\mathrm{KV}}d b n.\]

The factor two accounts for keys and values. Suppose \(L=32\),
\(H_{\mathrm{KV}}=8\), \(d=128\), and \(b=2\). Each token needs 131,072
bytes, or 128 KiB. An 8192-token sequence needs about 1 GiB. If the
\textbf{remaining KV budget after weights/workspaces} is 24 GiB, at most
24 such unshared sequences fit under this simplified bound. At 16,384
tokens, that drops to twelve.

Actual layouts may use sharding, quantization, shared prefixes, paging,
other attention architectures, and allocator overhead. Those require the
served model's measured configuration. PagedAttention addresses KV
memory management, not an unlimited-session promise. This arithmetic is
not a measurement of the installed Qwen model.
\href{https://arxiv.org/abs/2309.06180}{PagedAttention paper}.

\subsubsection{Diagnose a queue and a
bottleneck}\label{diagnose-a-queue-and-a-bottleneck}

In a stable system, Little's law gives average in-flight work
\(C=\lambda W\), where \(\lambda\) is completion rate and \(W\) is
average time in the measured system. Ten completed episodes per second
with forty seconds average residence implies 400 live episodes. Some may
be parked or evicted rather than retaining full KV; that choice changes
latency and recomputation. It does not disappear from accounting.

Suppose an environment executor supports four episodes per second, while
an inference service generates 400 output tokens per second and each
episode needs 500 output tokens. Ignoring prefill and other overhead
yields an inference ceiling of 0.8 episodes per second. Raising
environment throughput to eight cannot raise end-to-end throughput above
that ceiling. This is an upper bound, not a predicted benchmark.

Measure queued, admitted, prefill, decode, tool wait, grading,
persistence, and cleanup durations separately. Report warm/cold state,
input/output lengths, actual peak concurrency, and tail latencies. A
host semaphore cannot tell a serving replica how much a long session's
context will grow.

\subsubsection{Preserve the intended workload
mix}\label{preserve-the-intended-workload-mix}

Prefix affinity can improve reuse but may repeatedly favor short, common
tasks. Long-context sessions occupy KV longer and can starve. If the
trainer requested half simple and half complex groups, consuming 90\%
simple completions changes the training distribution even when raw
throughput rises. Admission needs capacity reservations, predicted
growth, queue age, and declared mixing rules. MISA-T studies
KV-residency-aware admission while preserving mixed workloads; its
paper's measured gains are not this project's results.
\href{https://arxiv.org/abs/2608.11152v1}{MISA-T}.

\textbf{Question.} A dashboard reports 99\% prefix-cache hits. Can
decode be faster by 99\%? \textbf{Solution.} No.~A hit indicates
reusable input state under a declared metric. Decode work, cache
lookup/eviction, memory contention, and output length remain. Track
saved prefill time and completed useful groups, not a hit percentage
alone.

\textbf{Deliverable.} A session-memory budget with overhead reserve,
admission and overload rules, phase timings, workload proportions, and
cancellation/recovery tests. The local concurrency and Docker receipts
do not establish a distributed GPU scheduler.

\subsection{Lab 8. Price a usable result with explicit
unknowns}\label{lab-8.-price-a-usable-result-with-explicit-unknowns}

Count all attempted work, including failures, rejected candidates,
setup, grading, and cleanup. A cost per successful clean trace needs a
denominator that defines useful output. A cheap unsafe episode is not a
saving. Model quality, reward calibration, and data rights can dominate
hardware cost.

For a declared hypothetical price sheet, write

\[C_{\mathrm{attempt}}=\frac{I_u p_u+I_c p_c+O p_o}{10^6}
+t_s p_s+C_{\mathrm{grade}}+C_{\mathrm{storage}}+C_{\mathrm{other}}.\]

\(I_u\) and \(I_c\) are uncached and billed-cached input tokens, \(O\)
is output tokens, and token prices are dollars per million. \(t_s p_s\)
prices sandbox residence in seconds. Add cache-write charges, minimum
billing intervals, retries, network transfer, and provisioned idle
capacity when applicable. Internal prefix reuse and provider invoice
categories are different facts. The equation is not any vendor's current
quote.

\subsubsection{Solve a per-attempt
example}\label{solve-a-per-attempt-example}

Assume twenty calls, each with 3000 uncached input and 200 output
tokens. Use invented prices of \$2 per million input and \$8 per million
output. Model cost is \(60{,}000\times2/10^6+4000\times8/10^6=\$0.152\).
Forty sandbox seconds at an invented \$0.001 per second add \$0.04. A
hypothetical \$0.01 grader charge gives \$0.202 per attempt, before
omitted costs.

If 20\% of attempts yield a clean, rights-approved, useful trace,
expected cost per such trace is \(0.202/0.2=\$1.01\) under steady
averages. Halving sandbox residence reduces that to
\(0.182/0.2=\$0.91\). Doubling useful yield to 40\% instead gives
\(0.202/0.4=\$0.505\), provided the accepted traces have equal quality
and no hidden selection cost. Better data yield can matter more than
faster setup.

If 90\% of input tokens receive a hypothetical cached price of \$0.50
per million, input cost becomes
\(6000\times2/10^6+54{,}000\times0.5/10^6=\$0.039\). That calculation
requires actual invoice-category counts and compatible cache conditions.
It cannot use a GPU prefix-hit chart as a substitute.

\subsubsection{Solve a break-even
example}\label{solve-a-break-even-example}

Suppose building and maintaining a simulator costs an invented \$12,000,
and two already-calibrated approaches produce equivalent usable outcomes
at \$1.50 and \$1.01 each. Ignoring future maintenance differences,
break-even is \(12{,}000/(1.50-1.01)\approx24{,}490\) usable outcomes.
If their fidelity, curation quality, failure risk, or ongoing costs
differ, the comparison is invalid until those differences enter the
model.

For self-hosted inference, allocated GPU-hour cost includes idle
capacity and failed work. Divide by completed accepted output over the
same interval. An already-owned local computer has no API invoice, but
electricity, occupancy, hardware, and operator time are not
mathematically zero. Do not equate zero marginal API charge with
commercial unit economics.

\subsubsection{Solve incomplete
accounting}\label{solve-incomplete-accounting}

Ten finalized calls report usage for only six. Adding those six counts
yields a known partial sum, not a complete cost. Without hard per-call
token caps and an applicable billing rule, missing counts have no
justified finite dollar upper bound. With caps, construct a conservative
interval and label its assumptions. A per-episode timeout is not a money
cap, particularly when provider retries or delayed cancellation can
still incur charges.

Keep prices, region, cache policy, billing unit, currency, collection
interval, useful-output definition, and omitted items in the receipt.
Use an explicit spend cap and authorized stop rule for any hosted
experiment. A trial budget authorizes neither more accounts nor a later
automatic rerun.

\textbf{Deliverable.} A spreadsheet-ready cost equation, three
sensitivity cases, equal-utility assumptions, missing-usage interval,
and a stop rule. Actual matched-study usage is explicitly incomplete, so
this workbook does not calculate its billed cost or a competitor
advantage.

\subsection{Lab 9. Review a deployment as several trust
boundaries}\label{lab-9.-review-a-deployment-as-several-trust-boundaries}

A threat model names assets, actors, attacker capabilities, trust
assumptions, controls, and residual risk. ``Runs in a container''
answers only part of that model. The policy, environment service,
grader, artifact collector, inference server, storage, and control API
have different permissions and failure modes.

\subsubsection{Produce a permission
matrix}\label{produce-a-permission-matrix}

{\def\LTcaptype{none} % do not increment counter
\begin{longtable}[]{@{}
  >{\raggedright\arraybackslash}p{(\linewidth - 4\tabcolsep) * \real{0.3333}}
  >{\raggedright\arraybackslash}p{(\linewidth - 4\tabcolsep) * \real{0.3333}}
  >{\raggedright\arraybackslash}p{(\linewidth - 4\tabcolsep) * \real{0.3333}}@{}}
\toprule\noalign{}
\begin{minipage}[b]{\linewidth}\raggedright
Component
\end{minipage} & \begin{minipage}[b]{\linewidth}\raggedright
Allowed data/operations
\end{minipage} & \begin{minipage}[b]{\linewidth}\raggedright
Must not receive
\end{minipage} \\
\midrule\noalign{}
\endhead
\bottomrule\noalign{}
\endlastfoot
Policy/agent & Public task, authorized observations, bounded tools &
Hidden tests, expected final state, control-plane credentials \\
Environment service & Own task state; enforce scoped transitions &
Cross-tenant state or unrestricted host commands \\
Grader & Trusted initial facts, detached state/history &
Agent-controlled scoring code or network authority \\
Artifact collector & Read frozen output, enforce complete archive bounds
& Answer keys; privileged mounts; writable agent continuation \\
Inference service & Authorized model context and serving weights &
Grader secrets or unrelated tenants' prompts \\
Controller & Admit/cancel jobs; own cleanup; persist receipts &
Unchecked user-selected image, mounts, or daemon flags \\
Operator API & Authorized tenant-scoped detail and exports & Anonymous
access to private snapshots or raw traces \\
\end{longtable}
}

Harbor documents separate verification as an opt-in mode; default shared
verification is not automatically isolated. Docker's security guidance
treats the daemon and configuration as privileged concerns. gVisor adds
a boundary against untrusted-code exploitation but documents broader
attack classes and limits.
\href{https://docs.harborframework.com/tasks/separate-verifier}{Harbor
verifier}, \href{https://docs.docker.com/engine/security/}{Docker
security},
\href{https://gvisor.dev/docs/architecture_guide/security/}{gVisor
security model}.

\subsubsection{Solve a boundary
exercise}\label{solve-a-boundary-exercise}

An agent cannot write the hidden grader file, but can write arbitrary
archive entries. It emits a symbolic link named \texttt{summary.csv},
plus 500 MiB under unrelated names. Reading only the selected filename
can follow the link or miss resource exhaustion. Validate the complete
detached archive before selecting output: reject links, traversal,
devices, duplicate/truncated entries, unsupported formats, and excess
total files/bytes. Freeze the agent first so checked files cannot change
between inspection and use.

A separate grader that still mounts the agent's writable workspace can
retain a race or import-path attack. A nonroot agent with the Docker
socket can retain host authority. A memory cgroup can limit resource use
without proving data isolation. Hashing a receipt proves byte
consistency relative to a known hash, not that a privileged host
produced honest facts.

Local synthetic Python worlds are trusted objects in one process; they
do not sandbox arbitrary Python. The existing Docker executor's
constraints and cleanup receipts are bounded evidence, not a
kernel-escape guarantee or multi-tenant service. Full operator snapshots
remain sensitive even when policy observations are partial and training
exports are disabled.

\subsubsection{Require evidence before
deployment}\label{require-evidence-before-deployment}

Use the following gates as an engineering review, not a certification
claim. NIST's control catalog provides a broader security/privacy
reference.
\href{https://csrc.nist.gov/pubs/sp/800/53/r5/upd1/final}{NIST SP
800-53}.

\begin{enumerate}
\def\labelenumi{\arabic{enumi}.}
\tightlist
\item
  Authenticate every control and export request; derive tenant scope
  from trusted identity rather than a client field.
\item
  Prove object-level authorization across database queries, artifacts,
  replay, and downloads. Test guessed IDs and mixed-tenant candidate
  pools.
\item
  Separate tenant credentials, filesystem roots, caches, logs, network
  policies, and retention. Review cache side channels and raw request
  leakage.
\item
  Pin reviewed images/dependencies, prevent arbitrary daemon parameters,
  and bound CPU, memory, PIDs, output, time, queued work, and spending.
\item
  Make cancellation, provider closure, persistence, and resource
  destruction observable. Preserve failures; release admission even
  after cleanup errors.
\item
  Exercise restart and crash recovery. Never resume an unknown live
  process or relaunch a paid attempt without a declared rule.
\item
  Complete the behavioral-fidelity study, independent grader red-team
  tests, expert outcome calibration, family-heldout evaluation, and
  rights review.
\item
  Enforce retention and deletion across source records, derived traces,
  exports, backups, and access logs. Deleting a row does not prove
  removal of its influence from an already-trained model.
\end{enumerate}

\textbf{Question.} Is a loopback-only API with no authentication safe to
expose through a reverse proxy? \textbf{Solution.} No.~Loopback location
is a deployment assumption, not user identity or authorization. Adding
TLS protects transport but does not prevent one authenticated tenant
from reading another's run.

\textbf{Deliverable.} A component permission matrix, attack tree,
residual-risk list, and readiness checklist whose items each link to a
test or operational proof. Mark missing evidence as blocked. A roadmap
entry, successful reference rollout, schema-valid trajectory, or
polished console cannot clear a deployment gate by itself.

\end{document}
